\subsection{Recursive least squares}

The discrete plant model in \eqref{eq:motor-discrete-transfer} gives the
one-step regression
\begin{equation}
  y(k)=C(k)x_e(k),
  \qquad
  C(k)=\begin{bmatrix}-y(k-1) & u(k-1)\end{bmatrix}.
  \label{eq:rls-regression}
\end{equation}
The recursive least-squares update with forgetting factor \(\alpha\) is
\begin{align}
  L(k) &=
  \frac{P(k-1)C^{T}(k)}
       {C(k)P(k-1)C^{T}(k)+\alpha},
  \label{eq:rls-gain}\\
  x_e(k) &= x_e(k-1)+L(k)\left(y(k)-C(k)x_e(k-1)\right),
  \label{eq:rls-parameter-update}\\
  P(k) &= \frac{\left(I-L(k)C(k)\right)P(k-1)}{\alpha}.
  \label{eq:rls-covariance-update}
\end{align}
The initial values are
\begin{equation}
  P(0)=100I,
  \qquad
  x_e(0)=\begin{bmatrix}-0.5&1.0\end{bmatrix}^{T}.
  \label{eq:rls-initialisation}
\end{equation}

Two implementation details are added to the textbook recursion. First, the
covariance update is suspended when the input/output data do not contain
excitation. Without this excitation gate, the division by \(\alpha<1\) in
\eqref{eq:rls-covariance-update} increases \(P\) without bound while the motor
is at rest. This covariance windup makes the subsequent parameter transient
too aggressive and can lead to bursting. Secondly, the controller-design step
uses a guard \(|b_0|>10^{-3}\). If the estimated numerator approaches zero, the
pole-placement formula would otherwise divide by a nearly singular plant gain.
